\chapter{Dữ Liệu và Pipeline Phân Tích}
\label{chap:data}

Chương này mô tả pipeline dữ liệu hiện hành từ một phòng Jitsi Meet đến Dashboard. Mục tiêu không chỉ là tạo nhãn sentiment, mà còn bảo toàn khả năng trả lời các câu hỏi: dữ liệu đến từ đâu, phiên bản nào đã được xử lý, vì sao một unit bị loại, adapter nào tạo prediction và kết quả nào còn hiệu lực. Cách tổ chức này xem provenance là thuộc tính cần bảo toàn của dữ liệu thay vì thông tin bổ sung sau suy luận \cite{buneman2001}.

\section{Nguồn dữ liệu, đơn vị xử lý và ranh giới trách nhiệm}
\subsection{Nguồn dữ liệu trong runtime}
Runtime hiện hành nhận URL Jitsi Meet trên miền \url{https://meet.jit.si/}, theo sau là tên phòng. Jitsi là lớp hội nghị WebRTC: tài liệu kiến trúc của Jitsi mô tả Jitsi Videobridge là media router, Jicofo điều phối conference và Prosody phục vụ signalling \cite{jitsiArchitecture2026}. Ứng dụng không triển khai lại các thành phần đó. Khi host và người tham gia đã được thông báo/đồng ý, backend yêu cầu bot Vexa tham gia phòng; bot chuyển audio window đến dịch vụ STT local và Vexa cung cấp snapshot transcript cho backend. Vexa có API/luồng self-host riêng cho bot và transcript \cite{vexaQuickstart2026}; trong đề tài, Vexa không phải database phân tích của ứng dụng.

Backend sao chép transcript vào PostgreSQL do ứng dụng sở hữu để version, phân tích lại và áp dụng retention/delete độc lập. REST reconcile là đường đồng bộ bắt buộc; signed webhook là tối ưu notification khi một endpoint HTTPS phù hợp đã được cấu hình. Cách tách này tránh giả định rằng một event của provider luôn đến đủ, đúng thứ tự hoặc không trùng.

Các dataset và notebook trong repository không được dùng làm ground truth cho chất lượng ứng dụng, vì domain, nhãn và đơn vị văn bản có thể khác transcript họp. Role-play hoặc transcript tổng hợp chỉ phù hợp với smoke test. Nếu xây dựng custom model, dataset phải ghi nhận license, mục đích, cách thu thập, thành phần và giới hạn theo tinh thần datasheets for datasets \cite{gebru2021}.

\subsection{Khái niệm và hợp đồng dữ liệu}
Pipeline dùng các khái niệm ở Bảng~\ref{tab:data-contract} để tách dữ liệu gốc, dữ liệu đủ điều kiện suy luận và dữ liệu trình bày. Hợp đồng thống nhất này là điểm nối của platform adapter: Jitsi là adapter runtime hiện hành; Google Meet, Teams và Zoom chỉ là hướng mở rộng sau khi có contract nguồn, consent và kiểm thử riêng.

\begin{table}[H]
\centering
\caption{Các thực thể và trường truy vết cốt lõi của pipeline}
\label{tab:data-contract}
\begin{tabularx}{\linewidth}{p{0.22\linewidth}X}
\toprule
Khái niệm & Định nghĩa trong hệ thống \\ \midrule
Raw segment & Segment transcript từ nguồn gồm định danh nguồn, speaker, text, timestamp, ngôn ngữ và trạng thái. Dữ liệu gốc được giữ để audit, không bị ``làm đẹp''. \\
Quality decision & Kết quả accepted/rejected cùng \texttt{quality\_reasons}; quyết định này tách khỏi raw text để nhìn thấy vì sao một đoạn không đi vào model. \\
Revision & Số phiên bản tăng khi text, speaker hoặc thời gian segment thay đổi. Meeting có \texttt{transcript\_revision} để gắn analysis run với ảnh chụp nguồn. \\
Analysis unit & Một đến ba final segment liền kề cùng speaker/ngôn ngữ, tối đa 45 giây; là đơn vị đầu vào của classifier. \\
Prediction & Label, confidence, phân bố lớp (nếu adapter cung cấp), polarity, inference time, model version và lý do abstain. \\
Evidence & Liên kết từ aggregate/prediction đến unit, segment ID, revision, text chuẩn hóa, timestamp, preprocessing version và model version. \\
Coverage & Tỷ lệ unit có prediction khác \textit{unknown}; không được ngầm diễn giải \textit{unknown} là neutral. \\
\bottomrule
\end{tabularx}
\end{table}

\section{Sơ đồ pipeline và quá trình vận hành}
Hình~\ref{fig:data-pipeline} biểu diễn đường dữ liệu hiện hành. PostgreSQL giữ bản ghi bền vững và outbox; Valkey Stream chỉ vận chuyển \texttt{run\_id}. Dashboard nhận tín hiệu SSE từ backend, sau đó đọc state có thẩm quyền qua REST thay vì gọi Vexa trực tiếp.

\begin{figure}[H]
\centering
\resizebox{\textwidth}{!}{%
\begin{tikzpicture}[
  node distance=9mm and 11mm,
  source/.style={draw=primaryblue, rounded corners=2pt, fill=lightblue, align=center, minimum width=26mm, minimum height=11mm},
  process/.style={draw=accentblue, rounded corners=2pt, fill=accentblue!12, align=center, minimum width=30mm, minimum height=11mm},
  store/.style={draw=primaryblue, cylinder, shape border rotate=90, aspect=0.28, fill=primaryblue!12, align=center, minimum width=25mm, minimum height=14mm},
  future/.style={draw=gray, dashed, rounded corners=2pt, fill=gray!10, align=center, minimum width=29mm, minimum height=11mm},
  flow/.style={-{Stealth[length=2.2mm]}, thick, draw=primaryblue},
  auxiliary/.style={-{Stealth[length=2.2mm]}, thick, draw=gray, dashed}
]
  \node[source] (meet) {Jitsi Meet\\host + participants};
  \node[source, right=of meet] (bot) {Vexa bot\\Jitsi adapter};
  \node[process, right=of bot] (stt) {STT local\\VAD + quality source gate};
  \node[process, right=of stt] (vexa) {Vexa transcript\\API};
  \node[process, right=of vexa] (api) {FastAPI\\sync / webhook};
  \node[source, above=of api] (dash) {Dashboard\\transcript, KPI, evidence};
  \node[store, below=of api] (db) {PostgreSQL\\segment, revision,\\run, outbox};
  \node[process, left=of db] (queue) {Valkey Stream\\analysis job};
  \node[process, left=of queue] (worker) {Sentiment worker\\quality + adapter};
  \node[future, above=of vexa] (future) {Google Meet / Teams / Zoom\\adapter (đề xuất)};

  \draw[flow] (meet) -- node[above, font=\scriptsize]{consent, audio} (bot);
  \draw[flow] (bot) -- (stt);
  \draw[flow] (stt) -- node[above, font=\scriptsize]{final segment} (vexa);
  \draw[flow] (vexa) -- (api);
  \draw[flow] (api) -- (db);
  \draw[flow] (db) -- (queue);
  \draw[flow] (queue) -- (worker);
  \draw[flow] (worker.south) -- ++(0,-6mm) -| node[pos=0.58, below, font=\scriptsize]{prediction, aggregate, outbox} (db.south);
  \draw[flow] (api) -- node[right, font=\scriptsize]{REST/SSE state} (dash);
  \draw[auxiliary] (future) -- (vexa);
\end{tikzpicture}%
}
\caption{Pipeline Jitsi--Vexa--STT--Meeting Intelligence. Nét đứt là adapter chưa được nghiệm thu trong runtime hiện hành.}
\label{fig:data-pipeline}
\end{figure}

Quy trình bắt đầu khi người dùng nhập URL Jitsi, chọn ngôn ngữ nói \texttt{vi}, \texttt{en} hoặc \texttt{auto}, và xác nhận consent. Backend kiểm URL, kiểm meeting active trùng phòng để tránh spawn lặp, sau đó gọi Vexa server-to-server. Khi transcript xuất hiện, segment được upsert theo khóa \texttt{(meeting\_id, vexa\_segment\_id)}. Thay đổi text, speaker hoặc timestamp tăng revision segment và \texttt{transcript\_revision} của meeting, đồng thời tạo outbox event có dedupe key.

Sau khi có segment final đạt điều kiện, backend tạo analysis run gắn với source revision, model version và preprocessing version, sau đó đưa \texttt{run\_id} vào Valkey Stream. Worker tạo analysis unit, gọi classifier adapter, lưu prediction và aggregate theo cửa sổ phút. Khi run thành công, worker phát event \texttt{analysis.sentiment.\allowbreak updated}; Dashboard nhận SSE như tín hiệu refresh và đọc lại overview, timeline và evidence qua API.

\section{STT local, quality gate và language routing}
\subsection{STT local và trade-off CPU}
Whisper cho thấy cách huấn luyện weak supervision quy mô lớn có thể hỗ trợ nhận dạng tiếng nói đa ngôn ngữ \cite{whisper2022}. Runtime này dùng \texttt{faster-whisper}, một triển khai Whisper trên CTranslate2. Tài liệu của dự án mô tả các tuỳ chọn CPU \texttt{int8}, language probability và VAD được dùng trong runtime \cite{fasterwhisper2026}.

Worker hiện dùng model \texttt{small}, CPU \texttt{int8}, một request active và VAD bảo thủ. Đây là lựa chọn vận hành để hạn chế queue backlog/503 trên máy CPU, không phải tuyên bố model này đạt chất lượng tương đương một dịch vụ transcription hoặc translation thương mại. Medium/large có thể cải thiện nhận dạng trong một số bối cảnh nhưng phải được benchmark trên phần cứng và workload thật trước khi thay runtime.

\subsection{Ba lớp lọc và bảo toàn audit}
Lớp STT bỏ audio ngắn, candidate có tín hiệu confidence xấu, decoded result quá ngắn và các mẫu media-outro đã biết như ``subscribe/đăng ký kênh''. Prompt của window trước bị tắt mặc định để một câu sai không quay lại làm context cho window sau. Mục tiêu là giảm khả năng một lỗi lặp biến thành transcript mới; đây không phải cơ chế chứng minh semantic correctness.

Backend vẫn giữ raw segment nhưng đặt cờ \texttt{analysis\_eligible} thành \texttt{false} nếu text rỗng, dưới ba token, chưa final, retracted, có duration không hợp lệ, chứa script không phù hợp, hoặc khớp mẫu hallucination. Dashboard ẩn các segment này mặc định nhưng cho phép người vận hành bật chúng để đối chiếu. Vì vậy, bộ lọc không âm thầm xóa bằng chứng nguồn và segment bị lọc không làm méo KPI/model.

Language router chuẩn hóa Unicode NFC/whitespace, xét dấu hiệu nội dung rồi đối chiếu metadata nguồn. Router tạo các route \texttt{vi}, \texttt{en}, \texttt{mixed} hoặc \texttt{unknown}. Baseline \texttt{baseline-rules-v3} nhận \texttt{vi}/\texttt{en}; route \texttt{mixed} và ngôn ngữ không hỗ trợ trả \textit{unknown} thay vì bị ép vào một nhãn polarity. Điều này phân biệt routing với dịch máy: Dashboard hiển thị transcript gốc, không tuyên bố có bộ dịch.

\subsection{Xây analysis unit, suy luận và tổng hợp}
Các segment accepted được ghép theo cùng speaker và language, tối đa ba segment, không vượt 45 giây. Khi speaker/language đổi, time window vượt ngưỡng hoặc gặp segment bị loại, unit đang mở được flush. Pipeline cũng khử duplicate gần nhau cùng speaker để nguồn phát lặp không làm méo aggregate. Những quy tắc này thuộc preprocessing version \texttt{2026-11-stt-quality-gate-v2}; thay đổi quy tắc phải tạo run mới thay vì ghi đè kết quả cũ.

Adapter hiện hành là \texttt{baseline-rules-v3}, một baseline theo luật Việt--Anh để kiểm chứng data contract, không phải mô hình đã huấn luyện. Custom model phải trả đúng hợp đồng prediction, có artifact hash, model/preprocessing version, tập đánh giá và giới hạn sử dụng được ghi nhận \cite{mitchell2019}. Prediction lưu theo cặp duy nhất \texttt{(analysis\_run\_id, analysis\_unit\_id)}; aggregate chỉ dùng nhãn \textit{positive}, \textit{neutral}, \textit{negative}. Nếu không có prediction đủ điều kiện, score là \texttt{null}, không phải neutral.

\section{Tính đúng đắn vận hành và giới hạn đo lường}
Pipeline hướng tới xử lý ít nhất một lần ở lớp queue. \texttt{analysis\_runs} có ràng buộc duy nhất trên bốn trường: định danh meeting, loại run, revision nguồn và version model; outbox có dedupe key. Worker chỉ acknowledge message sau khi transaction thành công. Nếu Valkey tạm không sẵn sàng khi enqueue, run vẫn được lưu queued để worker có thể khôi phục; điều này giải thích vì sao PostgreSQL, không phải queue, là nguồn dữ liệu chuẩn.

Độ trễ có thể phân rã thành:
\begin{equation}
  T_{e2e}=T_{capture}+T_{stt}+T_{ingest}+T_{queue}+T_{analysis}+T_{delivery}.
\end{equation}
Trong MVP chưa có benchmark tải hay log đầy đủ từng thành phần, nên chỉ có thể gọi hệ thống là \emph{near-real-time}. WER/CER chỉ được báo cáo khi có audio, transcript chuẩn và tokenization đã chốt; macro-F1/calibration chỉ được báo cáo với corpus sentiment đúng task và test set độc lập theo meeting. Không chỉ số nào trong các nhóm này được suy ra từ smoke test hoặc rule baseline.
